% Gold Standard GBFS France --- Patterns (Cell Press) submission draft
% Adapted from the SCITEPRESS conference version (gold_standard.tex)
% Single-column article class; Cell Press editors typeset the final layout.

\documentclass[11pt,a4paper]{article}

\usepackage[T1]{fontenc}
\usepackage[utf8]{inputenc}
\usepackage{lmodern}
\usepackage[margin=2.5cm]{geometry}
\usepackage{microtype}
\usepackage{amssymb,amsmath,amstext}
\usepackage{booktabs}
\usepackage{graphicx}
\graphicspath{{figures/}}
\usepackage{algorithm2e}
\usepackage{xcolor}
\usepackage{url}
\usepackage[colorlinks=true,linkcolor=blue,citecolor=blue,urlcolor=blue]{hyperref}
\usepackage{enumitem}
\usepackage{caption}
\usepackage{titlesec}
\usepackage{authblk}
\usepackage{mdframed}

\titleformat{\section}{\large\bfseries\sffamily\color{black!85}}{\thesection}{0.8em}{}
\titleformat{\subsection}{\normalsize\bfseries\sffamily\color{black!75}}{\thesubsection}{0.8em}{}

\newcommand{\ie}{\emph{i.e.}}
\newcommand{\eg}{\emph{e.g.}}
\newcommand{\etc}{etc.}

\newmdenv[
  linecolor=black!30,
  linewidth=0.6pt,
  backgroundcolor=black!3,
  innertopmargin=8pt,
  innerbottommargin=8pt,
  innerleftmargin=10pt,
  innerrightmargin=10pt,
  skipabove=10pt,
  skipbelow=10pt,
  frametitle={\textsf{\bfseries\large The Bigger Picture}},
  frametitlebackgroundcolor=black!8,
]{bigpicture}

\newenvironment{highlights}{%
  \par\noindent\textsf{\bfseries\large Highlights}\par\smallskip
  \begin{itemize}[leftmargin=1.4em,itemsep=0.15em,topsep=0.2em]}{%
  \end{itemize}\par\medskip}

\title{\bfseries\sffamily\Large Auditing Open Mobility Feeds at Country Scale:\\
A Reproducible Protocol and Reference Dataset for the 123 French GBFS Systems}

\author[1]{Rohan Foss\'e\thanks{Lead contact: \texttt{rfosse@cesi.fr}}}
\author[2]{Ga\"el Pallares}
\affil[1]{CESI Engineering School, Montpellier, France}
\affil[2]{CESI LINEACT (EA 7527), Montpellier, France}
\date{}

\begin{document}
\maketitle

\section*{Summary}
Open mobility feeds are increasingly assumed to be directly ingestible into
quantitative research pipelines: the General Bikeshare Feed Specification
(GBFS) is now a regulatory requirement for French operators, and 123 feeds
are centralised on a single public portal. By auditing all of them, we show
that this syntactic standardisation is necessary but not sufficient. Five
recurring anomaly classes (A1--A5: out-of-domain inclusion, placeholder
capacity, structural over-capacity, geospatial error, out-of-perimeter
coverage) re-appear across operators and together reclassify
\textbf{30.9\,\%} (95\,\% bootstrap CI [30.5, 31.3]) of the raw stations.
We formalise a six-step idempotent purging protocol, hybridise its output
with five reference sources (INSEE, BAAC, BD~TOPO, OpenStreetMap, FUB), and
release the resulting \emph{Gold Standard GBFS France} dataset (46{,}307
certified stations across 97 cities) as a versioned research commons under
ODbL on Zenodo. A single anomaly (the A3 reclassification of the
\emph{Pony} free-floating fleet in Bordeaux) is sufficient to move the city
from rank 2 to rank 14 in a national soft-mobility ranking, a result that
quantifies the cost of skipping the audit. We further publish three
internal robustness checks, a cross-country pilot on Bicing Barcelona
and Oslo Bysykkel (both pass the audit with zero anomalies under
French thresholds), a reproducibility checklist with container-pinned
bit-exact replay, and a six-experiment falsification roadmap with
four pilots already completed. The protocol is designed to
be transferable; the dataset itself is positioned as a contestable
artefact rather than a finished product.

\begin{bigpicture}\itshape
Public-policy decisions about urban mobility increasingly depend on the
analysis of open data feeds whose semantic quality is taken on trust.
This article makes the case that syntactic standardisation, on its own,
hides systematic errors that can flip the ranking of cities by an order
of magnitude, and proposes a reproducible audit pipeline backed by a
versioned reference dataset. The contribution is methodological as much
as empirical: any open mobility standard --- GBFS today, transit feeds,
charging-station feeds and shared-mobility-data-specification feeds
tomorrow --- benefits from a community-maintained validation tier
between regulatory publication and academic use. The 123-system French
corpus is a stress test of that idea, and the audit numbers documented
here support a broader conclusion: an open feed is a starting point for
research, not a substitute for one.
\end{bigpicture}

\begin{highlights}
\item Five anomaly classes reclassify 30.9\,\% of 67k raw GBFS stations
\item Six-step idempotent, reversible audit pipeline released as open code
\item Gold Standard dataset: 46{,}307 certified stations, 97 cities, FAIR
\item Bordeaux drops from national rank 2 to rank 14 after one correction
\end{highlights}

\noindent\textbf{Keywords:} data quality; GBFS; bike sharing; FAIR; reproducibility;
open data audit; reference dataset; mobility infrastructure.

% ===================================================================
\section{Introduction}\label{secIntro}

For more than a decade, bike-sharing systems have been a structural
component of urban mobility policy in France and across Europe, advancing
goals of liveability, decarbonisation and last-kilometre integration. To
support inter-operability of the resulting fleets, the General Bikeshare
Feed Specification (GBFS), introduced by the North American Bikeshare
Association in 2015 and now maintained by the Open Mobility Foundation,
provides a shared ontology covering station geolocation, dock counts,
real-time availability and pricing
\cite{Mobi2024,OMF2023}. In France, the 2019 Mobility Orientation Law
(LOM, article L.~1115-1 of the transport code) makes the publication of
these feeds on \emph{transport.data.gouv.fr} a regulatory acquis, so
that by 2026 the working assumption of many research and policy
pipelines is that GBFS feeds can be ingested directly.

That working assumption is misleading. The systematic audit reported
here, covering 123 French GBFS systems, shows that standardised syntax
masks substantial semantic heterogeneity. Three findings make the gap
concrete. First, the same field
(\texttt{station\_information.capacity}) covers radically different
realities: a physical dock count for dock-based fleets, a theoretical
estimator obtained by conditional averaging on free-floating fleets, and
arbitrary placeholders for under-reporting operators. Second, fourteen
\emph{Citiz} car-sharing systems are listed on the national portal under
the GBFS schema, blurring the very definition of the object of study.
Third, $3.8\,\%$ of the raw stations (95\,\% bootstrap CI [3.4, 4.2],
$n = 10{,}000$ resamples) exhibit transposed coordinates, fall outside
the national perimeter, or generate continental-scale bounding boxes.
Taken together, these observations support a single methodological
claim: syntactic standardisation of an open feed is not sufficient to
turn it into a research artefact, and any quantitative use of GBFS data
must be preceded by an explicit, documented and reproducible audit.

Empirical audits of mobility feeds have so far remained piecemeal. Tao
\textit{et al.} \cite{Tao2020} document sampling biases in mobile-phone
records used for origin-destination estimation; Romanillos \textit{et
al.} \cite{Romanillos2016} warn against statistical artefacts of
ingesting raw GPS feeds directly; Fishman \cite{Fishman2016} stresses
the heterogeneity of operational definitions in international BSS
comparisons. In the adjacent transit literature, the GTFS community has
been more proactive: the Canonical GTFS Schedule Validator
\cite{MobilityData2024Validator} and the earlier Conveyal validator
\cite{Antrim2013Conveyal} codify hundreds of schema-level checks on
transit feeds, and the data-engineering community has produced reusable
declarative-validation frameworks such as Great Expectations
\cite{GreatExpectations2024}, Pandera \cite{Pandera2024} and the
Frictionless Data specifications \cite{FrictionlessData2023};
ISO/IEC~25012 \cite{ISO25012} formalises the data-quality dimensions
that we implicitly mobilise. None of these references, however,
addresses GBFS specifically, and none of them produces a versioned,
audited reference dataset at the scale of a country. The present paper
fills that gap for the French case.

The contributions are fourfold. First, an empirically-grounded
\textbf{taxonomy} of five anomaly classes (A1--A5), derived from the
audit of all 123 French GBFS systems and stress-tested on two
non-French systems (Bicing Barcelona, Oslo Bysykkel) which both pass
the audit with zero anomalies under unchanged thresholds. Second, a
six-step \textbf{purging protocol} whose every transformation is
logged, reproducible and reversible. Third, the \textbf{Gold Standard
GBFS France} data product (46{,}307 certified stations, 97 cities),
released as a versioned Parquet file with DOI
\texttt{10.5281/zenodo.20125460} under ODbL, accompanied by its JSON
Schema, Croissant manifest and audit report. Fourth, a
\textbf{validation roadmap} of six falsifiable experiments, four of
which already include a completed pilot, so that the protocol can be
challenged in public rather than defended in private. Section~\ref{secResults} reports the empirical results,
Section~\ref{secDiscussion} discusses implications and limitations, and
Section~\ref{secMethods} details the experimental procedures.

% ===================================================================
\section{Results}\label{secResults}

\subsection{Anomaly taxonomy and incidence}
The systematic audit of the 123 French GBFS systems uncovered five
recurring anomaly classes, formalised in the \texttt{station\_type}
column of the Gold Standard and summarised in
Table~\ref{tab:anomalies}. Each class, in isolation, is sufficient to
invalidate a non-audited cross-system comparison; the protocol that
addresses them is given in Section~\ref{secMethods}.

\begin{table}[ht]
\centering
\caption{Five GBFS anomaly classes observed on the French corpus
(2025--2026). Incidence is reported at the system level except for A4
which is intrinsically station-level.}
\label{tab:anomalies}
\small
\begin{tabular}{@{}llll@{}}
\toprule
\textbf{Class} & \textbf{Name} & \textbf{Signature} & \textbf{Incidence}\\
\midrule
A1 & Out-of-domain inclusion & Car-sharing advertised as BSS & 14 systems\\
A2 & Placeholder capacity & Constant non-zero $c_i$, \eg $c = 100$ & 3 systems\\
A3 & Structural over-capacity & Conditional averaging on free-floating & 8 systems\\
A4 & Geospatial error & Transposed coords; $>3\sigma$ outliers; out-of-FR & $\sim 3.8\,\%$ of stations\\
A5 & Out of perimeter & System area $>50{,}000$\,km$^2$; overseas & 5 systems\\
\bottomrule
\end{tabular}
\end{table}

A3 is the most pernicious anomaly. Free-floating fleets (\emph{Pony},
\emph{Cykleo}, \etc) advertise virtual stations in GBFS, the term
``station'' designating the GPS position of an anchoring point
currently occupied by a vehicle. To avoid displaying empty stations,
the operator typically reports a capacity profile by \emph{conditional
averaging} on stations whose instantaneous capacity is non-zero:
\begin{equation}
\bar{c}_{\text{profile}}
= \frac{\sum_{i\,:\,c_i > 0} c_i}{\#\{i\,:\,c_i > 0\}}
\;\neq\;
\bar{c}_{\text{actual}}
= \frac{1}{N}\sum_{i=1}^{N} c_i.
\label{eq:a3-bias}
\end{equation}
Aggregated at the system level, $\bar{c}_{\text{profile}}$ wrongly
classifies thousands of free-floating bikes as heavy dock-based
stations. Cross-system comparison is therefore meaningful only after
explicit reclassification; the empirical signature is visible in
Figure~\ref{fig:capacity-violin} as a bimodal capacity distribution.

\begin{figure}[ht]
\centering
\includegraphics[width=0.75\linewidth]{fig02_capacity_violin.pdf}
\caption{Empirical signature of the A3 floating-anchor bias. Declared
capacity distribution by \texttt{station\_type} (clipped at 80 docks)
on the Gold Standard corpus. Only dock-based stations exhibit a
physically interpretable capacity (median = 20).}
\label{fig:capacity-violin}
\end{figure}

\subsection{Cumulative effect on the national corpus}
Table~\ref{tab:overview} summarises the cumulative effect of the
purging protocol. Across the country, 20{,}693 stations --- precisely
$30.9\,\%$ of the raw corpus (95\,\% bootstrap CI [30.5, 31.3],
$n = 10{,}000$ resamples on systems) --- are reassigned, removed or
reclassified. This is itself a strong empirical argument against the
view that the standard alone delivers a research-ready dataset.

\begin{table}[ht]
\centering
\caption{Effect of the purging pipeline on the national GBFS corpus.
\emph{Certified} entries are kept in the Parquet with a typed
\texttt{station\_type} so that downstream consumers can filter
explicitly; \emph{rejected} entries are kept in a separate
\texttt{rejected\_stations.parquet} with the exclusion motive.}
\label{tab:overview}
\small
\begin{tabular}{@{}lrr@{}}
\toprule
\textbf{Indicator} & \textbf{Raw} & \textbf{Gold Standard}\\
\midrule
Inventoried GBFS systems            & 142 & 123 (certified)\\
Cities covered                      & --  & 97 \\
\quad with at least one dock-based station & -- & 64 \\
Total stations (certified)          & 67{,}000 & 46{,}307 \\
\quad dock-based                    & --       & 5{,}442 \\
\quad free-floating (labelled, preserved) & -- & 39{,}235 \\
\quad car-sharing (A1)              & 1{,}630  & re-labelled, excluded from BSS use \\
Stations reclassified or dropped    & --       & 20{,}693 (30.9\,\%) \\
\bottomrule
\end{tabular}
\end{table}

The most heavily impacted urban areas are those operated by
free-floating fleets (\eg Bordeaux, Nice, Saint-\'Etienne) and those
historically lumped together with car-sharing systems on the national
portal (Strasbourg, Grenoble), whereas urban areas served by
traditional dock-based operators (V\'elib' in Paris, V\'eloMagg in
Montpellier) are left almost intact.

\subsection{Case study: Bordeaux and A3}
Bordeaux is the paradigmatic case (Table~\ref{tab:bordeaux}). The raw
feed assembles five operators across the agglomeration: \emph{Pony}
(free-floating), \emph{Bird} (free-floating), \emph{Dott}
(free-floating), \emph{Le V\'elo par TBM} (dock-based) and \emph{Citiz
Gironde} (car-sharing). \emph{Pony} alone advertises 2{,}996
\texttt{station\_information} entries with declared capacities around
12 docks. Naive ingestion therefore credits Bordeaux with tens of
thousands of nominal docking points and places it in the national top
three for V\'elib-class infrastructure.

\begin{table}[ht]
\centering
\caption{Bordeaux: impact of the A3 correction on the
\emph{Pony Bordeaux} subsystem and on the agglomeration as a whole.
The national IMD rank is reported for context (definition in
\cite{Fosse2026imd}).}
\label{tab:bordeaux}
\small
\begin{tabular}{@{}lrr@{}}
\toprule
\textbf{Indicator} & \textbf{Raw GBFS} & \textbf{Gold Standard}\\
\midrule
\textit{Pony Bordeaux subsystem} & & \\
\quad \texttt{station\_information} entries & 2{,}996 & 2{,}996 (relabelled) \\
\quad Declared as dock-based                 & 2{,}996 & 0 \\
\quad Declared capacity per entry (median)   & 12      & 12 \\
\quad Recomputed actual capacity             & --      & 0.03 bike/entry \\
\midrule
\textit{Bordeaux agglomeration (all 5 systems)} & & \\
\quad Total raw entries                      & 9{,}921 & -- \\
\quad Certified dock-based stations           & --      & 225 \\
\quad Nominal docking total (dock-based only) & 35{,}952 & 690 \\
\quad National IMD rank (companion study)     & 2       & 14 \\
\bottomrule
\end{tabular}
\end{table}

The purging protocol --- and in particular S3 (A3 correction) ---
reclassifies the 2{,}996 Pony entries as \texttt{free\_floating},
recomputes their actual capacity to 0.03 bike per entry, and leaves
the 225 genuinely dock-based stations of the historical \emph{V\'elo
par TBM} system. At the agglomeration level, this collapses the
nominal docking total by $98\,\%$, drops Bordeaux's dock-based
infrastructure from the top decile to the second quartile of the
national distribution, and shifts the Bordeaux IMD rank from 2 to 14
in our companion composite-index work \cite{Fosse2026imd}. The
twelve urban areas most exposed to A3 are mapped in
Figure~\ref{fig:bordeaux-before-after}.

\begin{figure}[ht]
\centering
\includegraphics[width=0.9\linewidth]{fig04_bordeaux_before_after.pdf}
\caption{Top twelve French urban areas ranked by absolute reduction in
dock-based station count between the raw GBFS feed and the certified
Gold Standard. Bordeaux exhibits the largest absolute drop
(9{,}921 raw entries down to 225 certified dock-based stations),
driven by the reclassification of the \emph{Pony} free-floating fleet
(A3). Paris and Lyon also see substantial corrections.}
\label{fig:bordeaux-before-after}
\end{figure}

\paragraph{Two complementary case studies.}
A3 is not the only class with a measurable downstream effect.
\emph{Nice} carries the flagship A2 placeholder via \texttt{pony\_Nice}:
every station declares $c = 100$, yielding a total nominal capacity of
about 25{,}000 bikes for an agglomeration whose actual dock-based
network is essentially the rebadged \emph{V\'eloBleu} of 91 stations.
After purging, the nominal capacity collapses by a factor of $250$ and
Nice's national rank moves out of the top quartile. \emph{Strasbourg}
provides the canonical A1 case: the \emph{Citiz Grand Est} car-sharing
system was historically ingested as part of the city's BSS inventory,
inflating the station count by roughly $35\,\%$. After semantic
exclusion, the agglomeration drops from rank 6 to rank 5 --- a smaller
absolute move than Bordeaux but on a more discriminating part of the
ranking. A1, A2 and A3 are therefore not interchangeable: each one
distorts the ranking through a different mechanism (false positive on
the population, on the capacity, or on the type), and a non-audited
pipeline conflates them.

\subsection{Comparison with naive baselines}
A natural question is whether the full pipeline is justified or whether
simpler heuristics would suffice. Three baselines are evaluated on the
same 123-system corpus
(Table~\ref{tab:baselines}). The \emph{Raw GBFS} baseline ingests
\texttt{station\_information} as is. The \emph{vehicle\_type filter}
mimics what a careful data engineer would write in a notebook. The
\emph{MobilityData audit} adapts the canonical GTFS-realtime
schema-level rules to GBFS.

\begin{table}[ht]
\centering
\caption{Measured comparison of the Gold Standard against three
baselines on the same 123-system French GBFS corpus. Brackets are
95\,\% confidence intervals (Wald on Dock, bootstrap on
$\rho_{\text{FUB}}$ over cities $n=2{,}000$, Wilson on A1/A3). ``Dock''
for the Gold Standard is the reference by construction.}
\label{tab:baselines}
\small
\begin{tabular}{@{}lcccc@{}}
\toprule
\textbf{Pipeline} & \textbf{Dock} & \textbf{$\rho_{\text{FUB}}$} &
\textbf{A1 ex.} & \textbf{A3 fix.}\\
\midrule
Raw GBFS &
  0.12 \footnotesize[.12, .12] &
  0.38 \footnotesize[.01, .68] &
  0\,\% & 0\,\% \\
\texttt{vehicle\_type} filter &
  0.15 \footnotesize[.15, .15] &
  0.19 \footnotesize[$-$.20, .52] &
  82\,\% \footnotesize[59, 94] & 0\,\% \\
MobilityData audit &
  0.15 \footnotesize[.15, .15] &
  0.19 \footnotesize[$-$.20, .52] &
  82\,\% \footnotesize[59, 94] & 0\,\% \\
\textbf{Gold Standard} &
  \textbf{1.00} &
  0.25 \footnotesize[$-$.13, .59] &
  \textbf{100\,\%} \footnotesize[82, 100] &
  \textbf{100\,\%} \footnotesize[91, 100] \\
\bottomrule
\end{tabular}
\end{table}

Three findings deserve emphasis. First, and most importantly, A3
reclassification recall is $0\,\%$ for every baseline (Wilson upper
bound $9\,\%$): structural over-capacity is invisible to syntactic
validators by construction. Only a semantically-aware step on the
capacity distribution can recover it. Second, the
\emph{vehicle\_type filter} and the \emph{MobilityData audit} baselines
produce identical numbers on this corpus, because schema-level checks
do not catch anything that the name filter did not already exclude ---
French GBFS feeds are syntactically well-formed but semantically
biased. Third, $\rho_{\text{FUB}}$ is highest on the raw count and
drops after audit; this is not a defect but a feature: FUB captures
perception of the entire cycling environment, so the raw count
over-fits the perceptual signal by including virtual stations, while
the Gold Standard's lower correlation reflects an honest count of
physical docks.

\subsection{Internal robustness checks}\label{secInternal}
Independent of the validation roadmap discussed in
Section~\ref{secDiscussion}, three internal robustness checks were
executed on the same corpus.

\paragraph{Orthogonal-coverage check.}
Three rule-based classifiers were implemented, each restricted to a
single signal source: a \emph{name-based} classifier seeing only
\texttt{system\_name} and \texttt{station\_name}; a
\emph{capacity-statistics} classifier seeing only the
\texttt{capacity} column; and a \emph{geospatial} classifier seeing
only $(\varphi,\lambda)$. The pairwise Cohen's $\kappa$ between the
three is near zero ($\kappa < 0.10$); Fleiss $\kappa = -0.07$. This is
not failure but the expected signature of complementary coverage: each
classifier sees only one slice of the taxonomy. The right measure is
the recall of their union against the integrated pipeline, which
reaches $97.2\,\%$ at a precision of $91.5\,\%$, with only 7 stations
out of 500 flagged by all three. This empirically demonstrates that
A1--A5 cannot be reduced to a single heuristic.

\paragraph{Full-task agreement check.}
Three independent decision-tree classifiers were implemented, each
seeing all signals but applying them in a different priority order
(name-first, capacity-first, geospatial-first). On the same 500-station
sample, the three classifiers reach Fleiss $\kappa = 0.80$, with
pairwise Cohen's $\kappa \in [0.70, 1.00]$. The 93/500 stations of
disagreement concentrate on the A1/A3 frontier (free-floating fleets
whose name does not match the floating-pattern dictionary). The lower
bound $\kappa = -0.07$ (orthogonal) and the upper bound $\kappa = 0.80$
(full-task) bracket the inter-annotator agreement that three independent
human annotators with full access to the payload would plausibly
achieve.

\paragraph{Threshold sensitivity.}
A single-dimension pilot of the topological threshold $\sigma_{\max}$
was completed on the certified corpus. Sweeping $\sigma_{\max} \in
\{2.0, 2.5, 3.0, 3.5, 4.0\}$ yields Jaccard similarities of
$0.66, 0.79, 1.00, 1.00, 1.00$ and Kendall's $\tau$ of
$0.92, 0.95, 1.00, 1.00, 1.00$ against the reference at
$\sigma_{\max} = 3.0$. Most importantly, the \textbf{Top-10 national
ranking is invariant} across the entire grid ($\sigma_{\max}$-induced
churn $= 0$), which is the criterion that matters for any downstream
public-policy reading.

These three checks do not replace human inter-rater reliability, which
is the next priority of the validation roadmap (experiment E1,
Section~\ref{secDiscussion}); they are reported as robustness
measures, not as substitutes.

\subsection{Cross-country transferability pilot}\label{secCrossCountry}
A natural objection to the present protocol is that A1--A5 might be
artefacts of French publication practice rather than structural
features of open-mobility-feed publication in general. To address this
upfront, two non-French dock-based systems were audited with the
\emph{same} rule set, recalibrating only the national perimeter
(latitude/longitude bounds). All other thresholds ($\sigma_{\max} = 3$,
$N_{\min} = 20$, A2 placeholder rule, A3 capacity-profile ratio
threshold of 5.0) were kept at their French values. Results are
reported in Table~\ref{tab:e5-pilot}; the full scripts and JSON audit
reports ship alongside the dataset under \texttt{experiments/e5\_europe/}.

\begin{table}[ht]
\centering
\caption{E5 pilot: applying the A1--A5 rule set to two non-French
dock-based systems, with French thresholds and country-appropriate
perimeter only. Both systems pass the audit with zero structural
anomalies, in contrast to the 22 French systems flagged by at least one
class. \emph{Ratio} is $\bar c_{\text{profile}} / \bar c_{\text{actual}}$
from Eq.~\ref{eq:a3-bias}: a value $\approx 1$ is the dock-based
signature, $> 5$ the floating-anchor signature.}
\label{tab:e5-pilot}
\small
\begin{tabular}{@{}lcccc@{}}
\toprule
\textbf{System} & \textbf{Country} & \textbf{Raw stns.}
                & \textbf{Capacity ratio} & \textbf{Anomalies}\\
\midrule
Bicing Barcelona       & ES & 544 & 1.006 & none of A1--A5 \\
Oslo Bysykkel          & NO & 266 & 1.000 & none of A1--A5 \\
\midrule
French corpus (ref.)   & FR & $\sim$67{,}000 & --- & 22 systems flagged \\
\bottomrule
\end{tabular}
\end{table}

Three observations follow. First, the capacity-profile ratio that
defines the A3 signature on free-floating French fleets is essentially
\textbf{unity} on both well-managed European dock-based systems, which
empirically validates the discriminative power of
Eq.~\ref{eq:a3-bias}. Second, vehicle-type metadata on both feeds is
populated cleanly (Bicing: \{\texttt{bicycle}, \texttt{scooter\_standing}\};
Oslo: \{\texttt{bicycle}\}), so no A1 false positive is triggered
either. Third, perimeter and topological filters caught only a
handful of edge outliers (1 station in Bicing, 3 in Oslo), in line
with what would be expected from any GPS-anchored fleet. These two
negative controls do not establish full European transferability ---
that requires the country-by-country E1 of the full E5 experiment ---
but they make it more difficult to argue that A1--A5 is over-fit to
the French market.

% ===================================================================
\section{Discussion}\label{secDiscussion}

\subsection{Implications}
The audit documented in this paper is intended as a scientific
validation tier, downstream of GBFS publication and upstream of
academic work. Its sustainability depends on collective appropriation,
since no single research team can durably operate a national
infrastructure on its own. Three concrete avenues are suggested.
First, a jointly published \emph{GBFS Audit Best Practices Handbook} by
francophone micromobility laboratories (CESI LINEACT, IFSTTAR
UMRESTTE, LAET, LVMT, IRSTEA-ITRT) would align practices currently
rebuilt in isolation. Second, a working group placed under ADEME or
\emph{transport.data.gouv.fr} could formalise an extended GBFS schema
and a certification framework, giving the audit a regulatory anchor.
Third, a continuous validation chain running the pipeline monthly on
the national feeds, with automated publication of the audit report and
operator-side notification on regression, would close the loop between
detection and correction. Star \cite{Star2010} reminds us that the
most durable infrastructures are those whose governance is explicit,
transparent and tied to an identifiable community of practice: moving
the Gold Standard from author-led to community-led governance is an
institutional challenge as much as a scientific one.

\subsection{Limitations}\label{secLimitations}
The protocol has four acknowledged limitations.

\textbf{L1 -- National coverage.} The corpus is limited to metropolitan
France; extension to overseas territories and other European
jurisdictions requires revisiting thresholds, in particular the
macro-regional surface threshold and the geofilter bounds.

\textbf{L2 -- Static versus dynamic.} The v1.0 release audits
\texttt{station\_information} only; auditing
\texttt{station\_status} (real-time availability) requires a distinct
protocol, sketched in experiment E4 (a 2-day pilot already detected
candidate dynamic anomalies on $12.9\,\%$ of monitored stations) and
developed in a companion paper \cite{Fosse2026flows}.

\textbf{L3 -- Free-floating asymmetry (key caveat).} The dock-based
subset (5{,}442 stations) is fully certified; the free-floating subset
(39{,}235 stations) is correctly typed but \emph{not rigorously
audited}. We have deliberately preserved free-floating entries in the
release with an explicit \texttt{station\_type = free\_floating} label,
so that downstream consumers can either exclude them or run their own
free-floating-native audit (experiment E6, scheduled for 2027). A
reader who interprets the Gold Standard as a fully-audited
all-fleet-types dataset will be misled; the appropriate framing is
\emph{Gold Standard for dock-based French BSS, with explicit
free-floating annotation}.

\textbf{L4 -- Operator coverage.} Some operators publish partial GBFS
feeds or are unavailable during the audit window; the inventory
completeness is therefore not guaranteed and must be re-evaluated
periodically, which is the purpose of the temporal-stability
experiment E3.

\subsection{Recommendations to operators}
The audit should also turn into constructive feedback to publishers.
Four recommendations follow directly from the taxonomy and are
deliberately compatible with the current GBFS specification: they are
conventions that the French operator community could adopt voluntarily,
without waiting for a schema revision.
\begin{enumerate}[leftmargin=*]
\item Populate \texttt{vehicle\_type} exhaustively, so that
car-sharing or scooter-sharing services on the portal can no longer
be ingested as bike sharing (A1).
\item Report \texttt{capacity} as a physical dock count for dock-based
stations and leave it \texttt{null} for virtual stations, rather than
filling it with placeholders such as $c = 100$ (A2).
\item Adopt the GBFS v3 \texttt{vehicle\_status} extension for
free-floating fleets, and stop publishing per-vehicle anchor positions
as \texttt{station\_information} (A3).
\item Publish geographical perimeters as polygons in the GBFS
\texttt{geofencing\_zones} block, so that validators can apply
system-specific topological filters (A4, A5).
\end{enumerate}

\subsection{Future directions}
Each of the four limitations above is paired with a falsifiable
experiment whose protocol, success criterion and (where available)
pilot results are documented and shipped alongside the dataset
(Table~\ref{tab:roadmap}). Four of the six experiments already
include a completed pilot (the internal robustness checks of
Section~\ref{secInternal} for E1, the $\sigma_{\max}$ sweep for E2,
the 2-day dynamic snapshot for E4, and the Bicing/Oslo cross-country
test of Section~\ref{secCrossCountry} for E5); the remaining
priorities are the human inter-rater reliability of E1, the full
four-dimensional threshold sweep of E2, the twelve-month temporal
stability of E3, and the free-floating-native protocol of E6. The
full protocols, including stratified sampling procedures, statistical
tests and success thresholds, are released under
\texttt{experiments/} in the source repository, so that the present
artefact can be challenged in public rather than defended in private.

\begin{table}[ht]
\centering
\caption{Six falsifiable experiments mapped to the limitations of
Section~\ref{secLimitations}. ``Pilot $\checkmark$'' indicates a
preliminary result already reported in this paper; full protocols are
documented under \texttt{experiments/} in the source repository.}
\label{tab:roadmap}
\small
\begin{tabular}{@{}llll@{}}
\toprule
\textbf{ID} & \textbf{Title} & \textbf{Targets} & \textbf{Status}\\
\midrule
E1 & Inter-rater reliability of A1--A5         & Auditor effect             & Internal pilot $\checkmark$ \\
E2 & Threshold and buffer sensitivity           & $\sigma_{\max}, N_{\min}$, radii & $\sigma_{\max}$ pilot $\checkmark$ \\
E3 & Twelve-month temporal stability            & Snapshot dependence (L4)    & Running \\
E4 & Dynamic extension to \texttt{station\_status} & L2                       & 2-day pilot $\checkmark$ \\
E5 & European generalisation (5 countries)      & L1                          & 2-system pilot $\checkmark$ \\
E6 & Free-floating-native protocol              & L3                          & Planned 2027-Q2 \\
\bottomrule
\end{tabular}
\end{table}

\subsection{Downstream use}
The Gold Standard is not produced for its own sake but to unlock
analyses that the raw feed would not support. Four families of
downstream studies illustrate the expected return on investment:
composite-index work (the companion IMD paper \cite{Fosse2026imd} is
the immediate example); demand modelling with graph-based neural
networks \cite{Lin2018BSSGraph} whose false-station nodes from A3 are
a known confound; spatial-equity analysis combining station access
with INSEE Filosofi; and multimodal integration with GTFS, since the
hybridisation schema is deliberately compatible with
\cite{MobilityData2024Validator}.

% ===================================================================
\section{Experimental Procedures}\label{secMethods}

\subsection{Provenance of the taxonomy}
Anticipating the auditor-effect critique, the iterative process that
produced A1--A5 is documented explicitly. The taxonomy was built by a
hybrid inductive-deductive workflow over five iterations between
2025-04 and 2026-01. \emph{Iteration 1} (April 2025): exploratory scan
of the 123 systems on the four obvious axes of ISO/IEC 25012
(completeness, consistency, accuracy, currentness); A1 and A4 emerge.
\emph{Iteration 2} (June 2025): enrichment with the system catalogue
of \emph{transport.data.gouv.fr}; A5 emerges as the system-level
complement of A4. \emph{Iteration 3} (September 2025): A2 discovered
when per-system standard deviation of \texttt{capacity} is plotted;
three systems with exact zero variance on non-zero capacities stand
out. \emph{Iteration 4} (November 2025): A3 emerges through reverse
engineering of free-floating fleets; the per-system median capacity of
\emph{Pony} matches a conditional-averaging estimator more closely
than a physical-dock count. \emph{Iteration 5} (January 2026): the
rule-based audit is re-run on a held-out subset of systems to test
class stability; two candidate classes do not survive (``A6-rolling
stock noise'', collapsed into A3, and ``A7-cross-border duplicate'',
no instance found). The current five-class taxonomy is the artefact of
that process; its provenance and rejected candidates ship with the
audit report.

\subsection{Six-step purging protocol}
The transition from the raw feed to the Gold Standard proceeds through
a chain of six sequential transformations (Algorithm~\ref{alg:purge}),
each coded in \texttt{utils/data\_loader.py} and logged in an audit
report shipped alongside the dataset. The protocol is deliberately
sequential: early steps are inexpensive and remove out-of-domain or
coarse-grained anomalies (A1, A2, A4 perimeter), while later steps
require system-level statistics ($\bar c$, centroids, $\sigma$) and
therefore benefit from operating on a pre-cleaned set.

\begin{algorithm}[ht]
\SetAlgoLined
\DontPrintSemicolon
\KwIn{National catalogue $\mathcal{C}$ (123 systems); thresholds
$\sigma_{\max} = 3$, $N_{\min} = 20$.}
\KwOut{Certified set $\mathcal{G}$, rejected set $\mathcal{R}$, audit report.}
$\mathcal{G} \gets \emptyset$; $\mathcal{R} \gets \emptyset$\;
\ForEach{system $s \in \mathcal{C}$}{
    Ingest GBFS feed $\to \mathcal{S}_s$\;
    \textbf{S1.} If $s$ is car-sharing, relabel \texttt{carsharing} and continue\;
    \textbf{S2.} If $\mathrm{Var}(c) = 0$ and $c > 0$, force \texttt{free\_floating}\;
    \textbf{S3.} Recompute $\bar c_{\text{actual}}$ on all stations and reassign topologically\;
    \textbf{S4.} Keep $\varphi_i \in [41^\circ,52^\circ]$, $\lambda_i \in [-6^\circ,10^\circ]$\;
    \textbf{S5.} Drop $d(\mathbf{x}_i, \boldsymbol{\mu}_s) > \sigma_{\max} \cdot \sigma_s$\;
    \textbf{S6.} If $|\mathcal{S}_s^{\text{dock}}| < N_{\min}$, label \texttt{too\_small}\;
    Update $\mathcal{G}, \mathcal{R}$ accordingly\;
}
\Return $(\mathcal{G}, \mathcal{R}, \text{audit report}, \text{metadata})$\;
\caption{GBFS-to-Gold-Standard purging pipeline.}
\label{alg:purge}
\end{algorithm}

\subsection{Protocol properties}
Four properties are designed into the pipeline. \emph{Idempotence:} two
successive applications produce the same result, tested by a unit test
executed on every CI run. \emph{Reversibility:} each excluded station
is kept in \texttt{rejected\_stations.parquet} together with the
exclusion motive. \emph{Logging:} every step records the number of
stations in, the number of stations out and the type of transformation
applied. \emph{Explicit parameterisation:} all thresholds
($\sigma_{\max}$, $N_{\min}$, geofilter bounds, hybridisation radii)
are documented as hyperparameters and can be varied for sensitivity
analysis without modifying the source code.

\subsection{Threshold justification}
The topological threshold $\sigma_{\max} = 3$ is a compromise between
robustness to GPS noise and the preservation of legitimate peripheral
extensions, \eg an off-centre TGV station or a university campus.
Section~\ref{secInternal} reports the single-dimension sensitivity
pilot: Top-10 invariance for $\sigma_{\max} \ge 2.5$. The robustness
threshold $N_{\min} = 20$ dock-based stations is calibrated on the
literature of meshed networks: below this value, density and
connectivity metrics become noise-dominated and are no longer
informative for cross-city comparison \cite{Saberi2018}.

\subsection{Multi-source hybridisation}
The output of the purging stage enriches each station with a
contextual profile obtained by articulating five reference sources
(Table~\ref{tab:sources}). Hybridisation is performed either by
spatial joining (point-in-polygon, buffer queries) or by administrative
aggregation (\ie at the IRIS or municipal level), depending on the
granularity of the source.

\begin{table}[ht]
\centering
\caption{Reference sources used for contextual enrichment.}
\label{tab:sources}
\small
\begin{tabular}{@{}llll@{}}
\toprule
\textbf{Source} & \textbf{Granularity} & \textbf{Variables} &
\textbf{Operation}\\
\midrule
INSEE Filosofi & IRIS / municipality &
   Median income/CU, poverty rate, local Gini & Aggregation \\
BAAC (Cerema)  & geolocated accident &
   Severe-crash density (500\,m buffer) & Spatial join \\
BD~TOPO        & polylines &
   Cycle-infrastructure linear (300\,m buffer) & Spatial join \\
OpenStreetMap  & polygons \& points &
   Facility density (transit stops, shops, services) & Spatial join \\
FUB Barometer  & municipal score &
   Subjective cycling perception (2017--2021) & Aggregation \\
\bottomrule
\end{tabular}
\end{table}

\subsection{Released schema}\label{secSchema}
The Gold Standard Parquet exposes 30 columns documented in
Table~\ref{tab:schema}: 14 from the raw GBFS payload (preserved
verbatim or normalised), 7 from the audit pipeline (typed,
\emph{station\_type} enum, audit status), 4 from spatial-equity
hybridisation, 3 from infrastructure hybridisation, and 2 from
multimodal-access hybridisation. Every column carries a stable name, a
declared dtype, a source pointer and a completeness rate; the
machine-readable equivalent is shipped as a JSON Schema, a Frictionless
Data Package descriptor and a Croissant manifest.

\begin{table}[ht]
\centering
\caption{Released schema of \texttt{stations\_gold\_standard\_final.parquet}.
``Compl.'' is the empirical completeness rate on the certified subset.
Columns marked $^\dag$ are introduced by the audit pipeline; columns
marked $^\ast$ are derived from external sources.}
\label{tab:schema}
\small
\begin{tabular}{@{}lllr@{}}
\toprule
\textbf{Column} & \textbf{Type} & \textbf{Source / Description} & \textbf{Compl.}\\
\midrule
\multicolumn{4}{l}{\textit{Identifiers (5)}}\\
\quad\texttt{uid}              & string  & Gold Standard primary key                 & 100\,\%\\
\quad\texttt{station\_id}       & string  & GBFS native identifier                    & 100\,\%\\
\quad\texttt{system\_id}        & string  & Operator-system identifier                & 100\,\%\\
\quad\texttt{system\_name}      & string  & Operator-system label                     & 100\,\%\\
\quad\texttt{source}            & string  & Feed URL / catalogue source               & 100\,\%\\
\midrule
\multicolumn{4}{l}{\textit{Spatial and administrative (5)}}\\
\quad\texttt{lat, lon}           & float64 & Geofiltered WGS84 coordinates             & 100\,\%\\
\quad\texttt{city}               & string  & Normalised city name                      & 100\,\%\\
\quad\texttt{commune\_name}      & string  & INSEE commune label                       & 99.1\,\%\\
\quad\texttt{code\_commune}      & string  & INSEE commune code                        & 99.1\,\%\\
\quad\texttt{region\_id}         & string  & Administrative region                     & 100\,\%\\
\midrule
\multicolumn{4}{l}{\textit{Station description (4)}}\\
\quad\texttt{station\_name}      & string  & GBFS station name                         & 99.8\,\%\\
\quad\texttt{address}            & string  & GBFS address                              & 41.2\,\%\\
\quad\texttt{capacity}           & int     & Declared capacity (may be a placeholder)   & 96.8\,\%\\
\quad\texttt{n\_stations\_system}& int     & Total stations in parent system            & 100\,\%\\
\midrule
\multicolumn{4}{l}{\textit{Audit outputs (2)$^\dag$}}\\
\quad\texttt{station\_type}      & enum    & \{docked\_bike, free\_floating, carsharing\} & 100\,\%\\
\quad\texttt{fetched\_at}        & datetime& Timestamp of the audited snapshot          & 100\,\%\\
\midrule
\multicolumn{4}{l}{\textit{Topography (2)$^\ast$}}\\
\quad\texttt{elevation\_m}                  & float64 & BD ALTI elevation                & 99.4\,\%\\
\quad\texttt{topography\_roughness\_index}  & float64 & Local relief amplitude           & 98.1\,\%\\
\midrule
\multicolumn{4}{l}{\textit{Cycling infrastructure (2)$^\ast$}}\\
\quad\texttt{infra\_cyclable\_km}           & float64 & BD TOPO cycle-lane linear (300\,m) & 97.6\,\%\\
\quad\texttt{infra\_cyclable\_pct}          & float64 & Share of dedicated right-of-way  & 96.9\,\%\\
\midrule
\multicolumn{4}{l}{\textit{Safety (1)$^\ast$}}\\
\quad\texttt{baac\_accidents\_cyclistes}    & int     & Severe-crash count (500\,m, 5\,y)  & 99.8\,\%\\
\midrule
\multicolumn{4}{l}{\textit{Multimodal access (2)$^\ast$}}\\
\quad\texttt{gtfs\_heavy\_stops\_300m}      & int     & Heavy-transit stops within 300\,m   & 98.7\,\%\\
\quad\texttt{gtfs\_stops\_within\_300m\_pct}& float64 & Share of accessible heavy-transit  & 98.7\,\%\\
\midrule
\multicolumn{4}{l}{\textit{Socio-economy (5)$^\ast$}}\\
\quad\texttt{revenu\_median\_uc}            & float64 & INSEE Filosofi median income/CU     & 95.4\,\%\\
\quad\texttt{gini\_revenu}                  & float64 & Local Gini index                    & 93.0\,\%\\
\quad\texttt{revenu\_d1}                    & float64 & First-decile income                 & 95.2\,\%\\
\quad\texttt{ecart\_interquar}              & float64 & Interquartile income spread         & 95.2\,\%\\
\quad\texttt{part\_menages\_voit0}          & float64 & Share of car-less households        & 95.4\,\%\\
\midrule
\multicolumn{4}{l}{\textit{Modal share (1)$^\ast$}}\\
\quad\texttt{part\_velo\_travail}           & float64 & Share of commute by bike (INSEE)    & 95.1\,\%\\
\bottomrule
\end{tabular}
\end{table}

\subsection{Reproducibility infrastructure}
The Gold Standard is released in alignment with the four FAIR
principles \cite{Wilkinson2016}. \emph{Findable:} primary deposit on
Zenodo under concept DOI \texttt{10.5281/zenodo.20125460}, with one
DOI per SemVer-tagged release (current 1.0.0, dated 2026-05); mirrored
on \emph{Recherche Data Gouv} (HAL/BSO indexing). \emph{Accessible:}
Parquet format, direct HTTPS download and Git LFS mirror.
\emph{Interoperable:} the schema extends GBFS v3 and is published as a
JSON Schema, a Frictionless Data Package descriptor
\cite{FrictionlessData2023}, a DCAT-AP record \cite{DCATAP2021} and a
Croissant manifest \cite{Croissant2024}. \emph{Reusable:} released
under ODbL v1.0; code under MIT, with audit report and source code.
Raw sources are archived with SHA-256 hashes; the Python environment
is pinned in \texttt{requirements.txt}; a Docker image is published
for bit-exact reproduction; random seeds are fixed (\texttt{numpy=42},
\texttt{python=42}); idempotence unit tests run in CI; the audit
report is generated automatically on every release.

\subsection{Engineering footprint}
The released Parquet weighs $7.8$\,MB on disk for the 46{,}307 stations;
it loads in $0.23$\,s into a \texttt{pandas} DataFrame with a memory
peak of $15.9$\,MB ($\sim 200{,}000$ rows/s on a commodity laptop).
Typical research queries are sub-25\,ms: filtering to dock-based
stations takes $5.2$\,ms, grouping by city $3.7$\,ms, a
Paris-bounding-box query $6.0$\,ms, a merge with FUB 2023 (62 cities)
$21.7$\,ms. The audit pipeline itself is implemented in $1{,}319$
lines of code across three modules (\texttt{utils/data\_loader.py},
\texttt{utils/gbfs\_collector.py}, \texttt{scripts/collect\_status.py}),
exposing 35 public functions and one collector class.

% ===================================================================
\section*{Resource Availability}

\paragraph{Lead contact.} Requests for further information should be
directed to and will be fulfilled by the lead contact, Rohan Foss\'e
(\texttt{rfosse@cesi.fr}).

\paragraph{Materials availability.} This study did not generate new
physical materials. The released artefact is the
\emph{Gold Standard GBFS France} dataset, a derivative of public open
data.

\paragraph{Data and code availability.}
\begin{itemize}[leftmargin=*]
\item \textbf{Dataset.} Deposited on Zenodo, concept DOI
\texttt{10.5281/zenodo.20125460}; current release 1.0.0 (2026-05),
licensed under ODbL v1.0. Mirrored on \emph{Recherche Data Gouv}.
\item \textbf{Source code.} The audit pipeline is available at
\url{https://github.com/rohanfosse/bikeshare-data-explorer}, licensed
under MIT, with a \texttt{CITATION.cff} file and a Docker image
(\texttt{ghcr.io/rohanfosse/gold-standard:1.0.0}) for bit-exact
reproduction. Releases are mirrored to Zenodo via the GitHub--Zenodo
integration.
\item \textbf{Schema and machine-readable descriptors.} JSON Schema,
DCAT-AP record, Frictionless Data Package descriptor and Croissant
JSON-LD manifest are shipped alongside the Parquet file.
\item \textbf{Raw inputs.} The 142 candidate GBFS feeds inventoried
from \emph{transport.data.gouv.fr} and MobilityData are archived with
SHA-256 hashes in the Zenodo deposit; the certified subset of 123
systems is derived deterministically from them.
\item \textbf{Additional resources.} The audit report, the
500-station blinded sample used in the internal robustness checks,
the orthogonal- and full-task classifier scripts and the
$\sigma_{\max}$ sensitivity notebook are released as supplementary
material alongside the dataset.
\end{itemize}

% ===================================================================
\section*{Usage Examples}\label{secUsage}

Three concrete reuse patterns are documented below; complete notebooks
are released alongside the dataset.

\paragraph{Example 1 -- Loading and filtering.}
Reading the certified Parquet and filtering to the dock-based subset
of a given city takes four lines:
\begin{quote}\small\ttfamily
import pandas as pd\\
gs = pd.read\_parquet("stations\_gold\_standard\_final.parquet")\\
paris = gs[(gs.city == "Paris") \&\\
\ \ (gs.station\_type == "docked\_bike")]\\
len(paris)\ \ \#\ returns\ 1{,}507
\end{quote}

\paragraph{Example 2 -- Spatial-equity analysis.} Per-city correlation
between dock-based station count and median income per consumption
unit:
\begin{quote}\small\ttfamily
by\_city = (\\
\ \ gs[gs.station\_type == "docked\_bike"]\\
\ \ .groupby("city").agg(\\
\ \ \ \ n=("uid", "size"),\\
\ \ \ \ income=("revenu\_median\_uc", "median"))\\
)\\
rho = by\_city.n.corr(\\
\ \ by\_city.income, method="spearman")\\
\#\ rho\ $\approx\ +0.39$\ (95\%\ bootstrap\ CI)
\end{quote}
A positive correlation of $\sim$0.4 indicates moderate income-density
coupling, consistent with the existence of \emph{social mobility
deserts} documented in the companion IMD paper \cite{Fosse2026imd}.

\paragraph{Example 3 -- Identifying mobility deserts.}
\begin{quote}\small\ttfamily
q1 = gs.revenu\_median\_uc.quantile(0.25)\\
deserts = gs[(\\
\ \ gs.station\_type == "docked\_bike") \&\\
\ \ (gs.revenu\_median\_uc < q1) \&\\
\ \ (gs.gtfs\_heavy\_stops\_300m == 0)]\\
deserts.city.value\_counts().head(10)
\end{quote}
This three-condition filter recovers $1{,}041$ stations
($19.1\,\%$ of the dock-based corpus) with a Q1 income threshold of
18{,}860\,EUR per consumption unit. Figure~\ref{fig:mobility-deserts}
shows the top fifteen contributing cities.

\begin{figure}[ht]
\centering
\includegraphics[width=0.75\linewidth]{fig06_mobility_deserts.pdf}
\caption{Top fifteen French urban areas by number of dock-based
mobility-desert stations, recovered by the three-condition filter of
Example 3. A station counts as a mobility desert if its commune sits
in the lowest national income quartile and the station has zero
heavy-transit stops within a 300\,m buffer.}
\label{fig:mobility-deserts}
\end{figure}

% ===================================================================
\section*{Acknowledgments}
The authors thank the producers of the open data that made this
research possible: GBFS operators, MobilityData, OpenStreetMap, Cerema,
ONISR, INSEE, FUB and the French GTFS community.

\section*{Author Contributions}
\textbf{R.F.}: Conceptualisation, Methodology, Software, Data curation,
Formal analysis, Writing -- original draft. \textbf{G.P.}:
Conceptualisation, Methodology, Supervision, Writing -- review \&
editing.

\section*{Declaration of Interests}
The authors declare no competing interests.

\bibliographystyle{unsrt}
\bibliography{references}

\end{document}
